\documentclass[10pt,a4paper]{amsart}
\usepackage{amsmath,amssymb,amsthm,graphicx}
\usepackage{verbatim}
\newenvironment{reference}{\comment}{\endcomment}
%\newenvironment{reference}{}{}
\newenvironment{slogan}{\comment}{\endcomment}
\newenvironment{history}{\comment}{\endcomment}

% For commutative diagrams we use Xy-pic
\usepackage[all]{xy}

% We use 2cell for 2-commutative diagrams.
\xyoption{2cell}
% Theorem environments.
%
\theoremstyle{plain}
\newtheorem{theorem}[subsection]{Theorem}
\newtheorem{proposition}[subsection]{Proposition}
\newtheorem{lemma}[subsection]{Lemma}

\theoremstyle{definition}
\newtheorem{definition}[subsection]{Definition}
\newtheorem{example}[subsection]{Example}
\newtheorem{exercise}[subsection]{Exercise}
\newtheorem{situation}[subsection]{Situation}

\theoremstyle{remark}
\newtheorem{remark}[subsection]{Remark}
\newtheorem{remarks}[subsection]{Remarks}

\numberwithin{equation}{subsection}

% Macros
%
\def\lim{\mathop{\mathrm{lim}}\nolimits}
\def\colim{\mathop{\mathrm{colim}}\nolimits}
\def\Spec{\mathop{\mathrm{Spec}}}
\def\Hom{\mathop{\mathrm{Hom}}\nolimits}
\def\Ext{\mathop{\mathrm{Ext}}\nolimits}
\def\SheafHom{\mathop{\mathcal{H}\!\mathit{om}}\nolimits}
\def\SheafExt{\mathop{\mathcal{E}\!\mathit{xt}}\nolimits}
\def\Sch{\mathit{Sch}}
\def\Mor{\mathop{\mathrm{Mor}}\nolimits}
\def\Ob{\mathop{\mathrm{Ob}}\nolimits}
\def\Sh{\mathop{\mathit{Sh}}\nolimits}
\def\NL{\mathop{N\!L}\nolimits}
\def\CH{\mathop{\mathrm{CH}}\nolimits}
\def\proetale{{pro\text{-}\acute{e}tale}}
\def\etale{{\acute{e}tale}}
\def\QCoh{\mathit{QCoh}}
\def\Ker{\mathop{\mathrm{Ker}}}
\def\Im{\mathop{\mathrm{Im}}}
\def\Coker{\mathop{\mathrm{Coker}}}
\def\Coim{\mathop{\mathrm{Coim}}}

% Boxtimes
%
\DeclareMathSymbol{\boxtimes}{\mathbin}{AMSa}{"02}

%
% Macros for moduli stacks/spaces
%
\def\QCohstack{\mathcal{QC}\!\mathit{oh}}
\def\Cohstack{\mathcal{C}\!\mathit{oh}}
\def\Spacesstack{\mathcal{S}\!\mathit{paces}}
\def\Quotfunctor{\mathrm{Quot}}
\def\Hilbfunctor{\mathrm{Hilb}}
\def\Curvesstack{\mathcal{C}\!\mathit{urves}}
\def\Polarizedstack{\mathcal{P}\!\mathit{olarized}}
\def\Complexesstack{\mathcal{C}\!\mathit{omplexes}}
% \Pic is the operator that assigns to X its picard group, usage \Pic(X)
% \Picardstack_{X/B} denotes the Picard stack of X over B
% \Picardfunctor_{X/B} denotes the Picard functor of X over B
\def\Pic{\mathop{\mathrm{Pic}}\nolimits}
\def\Picardstack{\mathcal{P}\!\mathit{ic}}
\def\Picardfunctor{\mathrm{Pic}}
\def\Deformationcategory{\mathcal{D}\!\mathit{ef}}
\usepackage{cleveref}
\crefname{originalreference}{original reference}{original references}
\Crefname{originalreference}{Original reference}{Original references}
\makeatletter
\newlabel{lemma-cone}{{08C5}{source}{}{sourcepointer-1}{}}
\newlabel{lemma-cone@cref}{{[originalreference][0][]08C5}{[1][0][]source}{}{}{}}
\makeatother
\crefname{originalreference}{original reference}{original references}
\Crefname{originalreference}{Original reference}{Original references}
\makeatletter
\newlabel{lemma-triseq}{{09QP}{source}{}{sourcepointer-2}{}}
\newlabel{lemma-triseq@cref}{{[originalreference][0][]09QP}{[1][0][]source}{}{}{}}
\makeatother
\crefname{originalreference}{original reference}{original references}
\Crefname{originalreference}{Original reference}{Original references}
\makeatletter
\newlabel{situation-ABC}{{09QJ}{source}{}{sourcepointer-3}{}}
\newlabel{situation-ABC@cref}{{[originalreference][0][]09QJ}{[1][0][]source}{}{}{}}
\makeatother
\title{A morphism of cone triangles whose first two components are homotopy equivalences has a homotopy equivalence on the cones}
\author{The Stacks Project Authors}
\date{stacks-project:dga.tex:local-snapshot-20261012; GNU FDL 1.2 or later}
\begin{document}
\maketitle
\noindent\textit{Editorial note.} All delivered mathematical text is the author's own, copied from the pinned Stacks Project chapter file sources/stacks-dga/pinned-original-source.tex; no statement or proof step was rewritten, reordered or removed.\par\medskip
\noindent\textit{Editorial note.} The counted proof is the author's complete argument: the inversion of a and b with the cone completion c' through lemma-cone, the reduction to the endomorphism case through lemma-triseq, and the difference trick with (c - 1)\textasciicircum{}2 = 0 giving the inverse 2 - c are all retained verbatim.\par\medskip
\noindent\textit{Editorial note.} The label lemma-analogue-third-isomorphism names the counted statement, and the proof environment immediately following it in the source is delivered verbatim as its proof.\par\medskip
\noindent\textit{Editorial note.} The counted statement is carried by the author's own lemma environment in the chapter source.\par\medskip
\noindent\textit{Allowed standard prerequisites (editorial).}\begin{itemize}
\item Completing a morphism on cones. Same chapter, Lemma lemma-cone. Builds the inverse component c' on the cones from the inverse homotopy equivalences a, b, through the left commutative box.
\item The two-out-of-three reduction for triangle morphisms. Same chapter, Lemma lemma-triseq. Reduces the assertion to the endomorphism case (1, 1, c) and yields the identity (c - 1)\textasciicircum{}2 = 0 from the admissible rows of the difference diagram.
\end{itemize}\medskip
\ifdefined\hypertarget\hypertarget{sourcepointer-1}{}\fi
\noindent\textit{Source reference.} Original reference 08C5: sources/stacks-dga/pinned-original-source.tex, source lines 3779--3779. Prerequisite pointer invoked at proof line(s) 4285; it names a standard fact of the same chapter that the counted argument uses, and every step of the counted proof built on it is retained verbatim in the delivered body.\par\medskip
\ifdefined\hypertarget\hypertarget{sourcepointer-2}{}\fi
\noindent\textit{Source reference.} Original reference 09QP: sources/stacks-dga/pinned-original-source.tex, source lines 4050--4050. Prerequisite pointer invoked at proof line(s) 4287; it names a standard fact of the same chapter that the counted argument uses, and every step of the counted proof built on it is retained verbatim in the delivered body.\par\medskip
\ifdefined\hypertarget\hypertarget{sourcepointer-3}{}\fi
\noindent\textit{Source reference.} Original reference 09QJ: sources/stacks-dga/pinned-original-source.tex, source lines 3762--3762. Prerequisite pointer invoked at proof line(s) 4262; it names a standard fact of the same chapter that the counted argument uses, and every step of the counted proof built on it is retained verbatim in the delivered body.\par\medskip
\begin{lemma}
\label{lemma-analogue-third-isomorphism}
In Situation \ref{situation-ABC} let $f_1 : x_1 \to y_1$ and
$f_2 : x_2 \to y_2$ be morphisms in $\text{Comp}(\mathcal{A})$. Let 
$$
(a,b,c): (x_1,y_1,c(f_1), f_1, i_1, p_1) \to (x_2,y_2, c(f_2), f_2, i_1, p_1)
$$
be any morphism of triangles in $K(\mathcal{A})$.
If $a$ and $b$ are homotopy equivalences, then so is $c$. 
\end{lemma}

\begin{proof}
Since $a$ and $b$ are homotopy equivalences, they are invertible in
$K(\mathcal{A})$ so let $a^{-1}$ and $b^{-1}$ denote their inverses
in $K(\mathcal{A})$, giving us a commutative diagram 
$$
\xymatrix{
x_2\ar[d]^{a^{-1}}\ar[r]^{f_2} &
y_2\ar[d]^{b^{-1}}\ar[r]^{i_2} &
c(f_2)\ar[d]^{c'} \\
x_1\ar[r]^{f_1} &
y_1 \ar[r]^{i_1} &
c(f_1)
}
$$
where the map $c'$ is defined via Lemma \ref{lemma-cone} applied to the left
commutative box of the above diagram. Since the diagram commutes
in $K(\mathcal{A})$, it suffices by Lemma \ref{lemma-triseq} to
prove the following: given a morphism of triangle
$(1,1,c): (x,y,c(f),f,i,p)\to (x,y,c(f),f,i,p)$
in $K(\mathcal{A})$, the map $c$ is an isomorphism in
$K(\mathcal{A})$. We have the commutative diagrams in $K(\mathcal{A})$:
$$
\vcenter{
\xymatrix{
y\ar[d]^{1}\ar[r] &
c(f)\ar[d]^{c}\ar[r] &
x[1]\ar[d]^{1} \\
y\ar[r] &
c(f) \ar[r] &
x[1]
}
}
\quad\Rightarrow\quad
\vcenter{
\xymatrix{
y\ar[d]^{0}\ar[r] &
c(f)\ar[d]^{c-1}\ar[r] &
x[1]\ar[d]^{0} \\
y\ar[r] &
c(f) \ar[r] &
x[1]
}
}
$$
Since the rows are admissible short exact sequences, we obtain
the identity $(c-1)^2 = 0$ by Lemma \ref{lemma-triseq}, from
which we conclude that $2-c$ is inverse to $c$ in $K(\mathcal{A})$
so that $c$ is an isomorphism. 
\end{proof}

\end{document}
